% !TEX TS-program = pdflatex
% CV ciblé — Ingénieur Infrastructure — Tunisie Leasing et Factoring
\documentclass[a4paper,9pt]{article}
\usepackage[utf8]{inputenc}
\usepackage[T1]{fontenc}
\usepackage[scaled=0.90]{helvet}
\renewcommand{\familydefault}{\sfdefault}
\usepackage[left=0.8cm,right=0.8cm,top=0.3cm,bottom=0.25cm]{geometry}
\usepackage{enumitem}
\usepackage{titlesec}
\usepackage{xcolor}
\usepackage[hidelinks]{hyperref}
\definecolor{navy}{RGB}{23,54,93}
\pagestyle{empty}
\setlength{\parindent}{0pt}
\setlength{\parskip}{0pt}
\linespread{0.82}
\hypersetup{pdftitle={CV - Baha Eddine Belhaj Mouldi - Ingenieur Infrastructure - Tunisie Leasing et Factoring},pdfauthor={Baha Eddine Belhaj Mouldi},pdfsubject={Infrastructure, Systemes, Reseaux, Virtualisation, Securite, Supervision}}
\titleformat{\section}{\normalsize\bfseries\color{navy}}{}{0pt}{\MakeUppercase}[{\vspace{1pt}\color{navy}\hrule height 0.6pt}]
\titlespacing*{\section}{0pt}{1.5pt}{0.8pt}
\setlist[itemize]{leftmargin=1.0em,label=\textbullet,itemsep=0pt,topsep=0pt,parsep=0pt}
\newcommand{\cventry}[4]{\textbf{#1} \hfill \textbf{#4}\\[0.2pt]#2{\ifx\relax#3\relax\else, #3\fi}\par\vspace{0.3pt}}
\begin{document}
\begin{center}
{\LARGE\bfseries\color{navy} Baha Eddine Belhaj Mouldi}\\[2pt]
{\normalsize\color{navy} Ingénieur Infrastructure \textbar\ Systèmes \textbar\ Réseaux \textbar\ Virtualisation \textbar\ Sécurité \& Supervision}\\[3pt]
{\small Tunis, Tunisie \quad|\quad +216 55 128 982 \quad|\quad \href{mailto:bahaeddine.belhajmouldi@gmail.com}{bahaeddine.belhajmouldi@gmail.com}}\\[1pt]
{\small \href{https://www.linkedin.com/in/bahamouldi}{linkedin.com/in/bahamouldi} \quad|\quad \href{https://github.com/bahamouldi}{github.com/bahamouldi} \quad|\quad \href{https://bahamouldi.github.io/portfolio/}{bahamouldi.github.io/portfolio}}
\end{center}
\vspace{1pt}
\section{Profil}
Ingénieur en génie informatique (ENIT, mention Très Bien) avec une expérience concrète de conception, déploiement et administration d'infrastructures systèmes et réseaux. Pratique de l'administration Windows Server / Linux, des environnements virtualisés, de la gestion des sauvegardes et de la continuité d'activité, acquise en travaux pratiques puis approfondie sur le terrain pendant un an de mission freelance et durant le stage de fin d'études chez Beehive Entreprises. Expérience directe de la haute disponibilité et de la reprise après sinistre via le déploiement en production d'un cluster Kubernetes zéro downtime, ainsi que de la cybersécurité et de la gestion des correctifs (37 CVE patchés) au travers du projet BeeWAF, conforme à 7 référentiels de sécurité. Esprit analytique, rigueur, autonomie et goût du travail en équipe.
\section{Compétences techniques}
\textbf{Systèmes} : Installation et administration Windows Server, Active Directory (utilisateurs, groupes, GPO), Linux (Ubuntu, CentOS, Debian, Kali), durcissement système\\[0.3pt]
\textbf{Infrastructures physiques \& virtualisées} : environnements virtualisés (Hyper-V --- notions), Docker, Kubernetes (Deployment, Service, Ingress, PVC, Secrets)\\[0.3pt]
\textbf{Haute disponibilité \& reprise après sinistre} : déploiement zéro downtime (maxUnavailable: 0, maxSurge: 1), health/liveness probes, dimensionnement des ressources, cluster multi-nœuds\\[0.3pt]
\textbf{Sauvegarde \& continuité} : solutions de sauvegarde et restauration (Veeam), plans de continuité d'activité\\[0.3pt]
\textbf{Réseaux} : TCP/IP, DNS, DHCP, VPN, Proxy, Firewall, segmentation réseau, Active Directory\\[0.3pt]
\textbf{Sécurité \& correctifs} : gestion des accès utilisateurs, politiques de sécurité, cybersécurité, virtual patching de CVE, PKI, MFA, IAM, conformité ISO 27001/NIST/PCI DSS/GDPR/SOC 2\\[0.3pt]
\textbf{Supervision \& performance} : Prometheus, Grafana, ELK Stack (Elasticsearch, Logstash, Kibana), monitoring temps réel, optimisation des performances serveur\\[0.3pt]
\textbf{Scripting} : Python, Bash, PowerShell, automatisation d'exploitation
\section{Expérience professionnelle}
\cventry{Ingénieur Sécurité \& Infrastructure --- Stagiaire PFE (mention Très Bien)}{Beehive Entreprises}{Tunisie}{02/2026 -- 06/2026}
\begin{itemize}
  \item Conception et déploiement d'une infrastructure en production sur cluster Kubernetes (5 nœuds) : haute disponibilité, zéro downtime, dimensionnement des ressources, health probes.
  \item Supervision et optimisation des performances : monitoring Prometheus/Grafana + ELK Stack, alerting temps réel, maintenance préventive.
  \item Gestion des correctifs et politiques de sécurité : virtual patching de 37 CVE, conformité continue sur 7 référentiels (OWASP, PCI DSS, GDPR, NIST, ISO 27001, CIS, SOC 2).
\end{itemize}
\cventry{Spécialiste Infrastructure \& Sécurité, Freelance (1 an)}{Beehive Entreprises}{Tunisie}{01/2025 -- 12/2025}
\begin{itemize}
  \item Installation et administration Windows Server / Active Directory (utilisateurs, groupes, GPO) et Linux ; configuration réseau DHCP, DNS, VPN, proxy et règles de pare-feu pour des infrastructures clients.
  \item Mise en œuvre de solutions de virtualisation et de sauvegarde (Veeam) ; gestion des accès utilisateurs et audit de vulnérabilités sur 8 applications, rapports de remédiation transmis aux équipes techniques.
\end{itemize}
\cventry{Stagiaire Cybersécurité --- Détection \& Réponse à incident SAP}{Streamlink}{Tunisie}{07/2025 -- 08/2025}
\begin{itemize}
  \item Supervision automatisée (Python/PyRFC, ELK, N8N) de transactions SAP critiques : détection $<$2 min, blocage automatisé $<$7 min.
\end{itemize}
\cventry{Chef de projet technique, Indépendant --- Équipe de 8 personnes}{VMate}{Tunisie}{01/2024 -- 06/2024}
\begin{itemize}
  \item Coordination d'une équipe pluridisciplinaire de 8 personnes : planification, arbitrages techniques, reporting --- travail en équipe et communication technique au quotidien.
\end{itemize}
\cventry{Stagiaire Sécurité Réseau}{Tunisie Telecom}{Tunisie}{07/2024}
\begin{itemize}
  \item Analyse de trafic réseau, durcissement de configuration, identification de 10+ vulnérabilités, dashboards de supervision Elastic/Kibana.
\end{itemize}
\section{Projets \& travaux pratiques}
\cventry{Infrastructure Windows Server / Active Directory \& Réseau Périmétrique}{Windows Server, Active Directory, GPO, DHCP, DNS, VPN, Proxy, Firewall}{}{ENIT}
\begin{itemize}
  \item Administration Active Directory (utilisateurs, groupes, unités d'organisation, GPO), configuration DHCP/DNS et sécurisation périmétrique (VPN, proxy, règles firewall, segmentation réseau).
\end{itemize}
\cventry{BeeWAF --- Infrastructure de Sécurité Haute Disponibilité}{FastAPI, Docker, Kubernetes, Jenkins, ArgoCD, Prometheus, Grafana}{}{2026}
\begin{itemize}
  \item Cluster Kubernetes en production, zéro downtime, supervision temps réel, gestion des correctifs (37 CVE), conformité continue sur 7 référentiels réglementaires.
\end{itemize}
\section{Formation}
\cventry{Diplôme d'Ingénieur en Génie Informatique}{École Nationale d'Ingénieurs de Tunis (ENIT)}{Tunisie}{09/2023 -- 06/2026}
\begin{itemize}
  \item Diplômé le 25/06/2026, mention Très Bien. Classé 65e sur 600+ au concours national d'entrée. Modules \& TP : systèmes et réseaux (Windows Server, Active Directory, DHCP, DNS, VPN, firewall), virtualisation, cryptographie et PKI, sécurité des systèmes d'information.
\end{itemize}
\cventry{Classes préparatoires Physique-Technologie}{Institut Préparatoire aux Études d'Ingénieurs, El Manar}{Tunisie}{09/2021 -- 06/2023}
\cventry{Baccalauréat Technique --- Mention Très Bien}{Lycée Abdelaziz Khouja, Kélibia}{Tunisie}{06/2021}
\section{Certifications}
Réseaux (CCNA) --- Cisco (2025) ; Fondements du Deep Learning --- NVIDIA (2025) ; IA Adversariale --- NVIDIA (2025) ; Machine Learning --- NVIDIA (2024) ; Git \& GitHub --- 365 Data Science (2024).
\section{Langues}
Arabe (Natif) \quad|\quad Français (Courant) \quad|\quad Anglais (B2)
\end{document}
